\chapter*{Abstract}

Large language models can now generate complete backend applications, but
existing benchmarks evaluate them only in isolated, single-container
environments and stop at functional and security correctness. Production
backends run under shared cluster resources, where replica counts,
resource limits, pod placement, and database topology jointly determine
achievable throughput: decisions usually made manually or by reactive
autoscalers unaware of the interaction between them. This thesis asks
whether an LLM agent can close that gap: iteratively optimise deployment
configuration on a Kubernetes cluster, guided by runtime performance
feedback, to maximise sustainable throughput under fixed infrastructure
constraints.

We extend \textsc{BaxBench} with a Kubernetes-based iterative
benchmarking framework in which an agent alternates, one lever per
iteration, between refining application code and a constrained deployment
specification, guided by load-test and resource feedback after each
deploy-and-benchmark cycle. Performance is scored by \emph{sustained
goodput}, discovered automatically per deployment without per-scenario
tuning.

We evaluate three frontier-tier models (Claude Opus 4.8, GPT-5.5, GLM-5.2)
across eight \textsc{BaxBench} scenarios and three application frameworks,
an evaluated matrix of up to 72 scenario$\times$framework$\times$model
cells and ten refinement iterations each, 63 of which produced experiment
data. Iterative refinement improves sustained goodput in essentially
every cell with data, in several cases converting a non-functional
deployment into a working one; gains range from just over $1\times$ to
more than $75\times$ depending on how under-provisioned the baseline was.
Code refinement contributes the largest single-step recoveries, dominated
by correctness bugs invisible to single-container testing; deployment
refinement is a smaller capacity knob that tracks workload characteristics
without being given target ranges, though its net effect varies by model,
turning net negative for one of the three. All three models reach a
working baseline and complete their full refinement budget on every cell
with data; the nine cells without data are concentrated almost entirely in
a single scenario, leaving the other seven fully usable for comparison.
LLM-guided deployment optimisation is a viable, practically bounded
complement to code-level optimisation, whose largest current limitation is
entanglement with code correctness rather than the deployment-reasoning
task itself.
